\apartado{2.7}{Variables aleatorias}{distrib}
\bajada{Hasta aquí preguntábamos por eventos sueltos. Ahora describimos el
comportamiento completo de una variable de una sola vez, y lo resumimos en dos
números que vas a ver en cada artículo que leas: la media y la desviación.}

\definicion{Variable aleatoria discreta}{%
  Una regla que asigna un número a cada resultado del experimento, y que toma
  valores separados: se cuentan, no se miden. El puntaje de un cuestionario, el
  número de sesiones a las que asiste un paciente, la cantidad de casos que
  llegan en una semana.}

\minihistoria{Dos números que resumen una distribución entera}{%
  Cuando un artículo dice «los participantes obtuvieron $M=6{,}32$
  ($DE=4{,}64$)», está resumiendo doscientos números en dos.\par
  El primero dice alrededor de qué valor se agrupan; el segundo, cuánto se
  dispersan. Con esos dos, y suponiendo una forma de campana, se puede
  reconstruir casi toda la distribución — que es lo que hace posible el
  apartado 2.9.}

\ejercicio{0}{¿Cuáles de estas son variables aleatorias \emph{discretas} y
  cuáles \emph{continuas}? Explica en una línea cada decisión.
  \begin{enumerate}[label=(\alph*),topsep=2pt,itemsep=1pt]
    \item El puntaje de un cuestionario de 9 ítems (0 a 27).
    \item El tiempo que tarda una persona en completarlo.
    \item La cantidad de sesiones a las que asiste un paciente.
    \item El nivel de cortisol en saliva.
  \end{enumerate}}{}
\renglones[4]
\respuesta{Discretas: (a) y (c). Continuas: (b) y (d).}
\solucion{%
  (a) \textbf{Discreta}: sólo puede valer $0,1,\ldots,27$. No existe el puntaje
  $14{,}3$.\par
  (b) \textbf{Continua}: entre 4 y 5 minutos hay infinitos valores posibles.\par
  (c) \textbf{Discreta}: se cuentan sesiones, y no hay media sesión.\par
  (d) \textbf{Continua}: es una medición física, limitada sólo por la precisión
  del instrumento.
  \par
  La distinción importa porque cambia la herramienta: las discretas usan
  binomial, Poisson o hipergeométrica (2.8); las continuas, la normal (2.9).}

\ejercicio{1}{Un programa breve tiene 4 sesiones. Sea $X$ el número de sesiones
  a las que asiste un paciente, con esta distribución:
  \begin{center}\small
  \begin{tabular}{@{}lccccc@{}}
  \toprule
  $x$ & 0 & 1 & 2 & 3 & 4\\
  $P(x)$ & $0{,}10$ & $0{,}15$ & $0{,}20$ & $0{,}25$ & $0{,}30$\\
  \bottomrule
  \end{tabular}
  \end{center}
  \textbf{Antes de calcular, escribe a cuántas sesiones crees que asiste un
  paciente en promedio:} \underline{\hspace{1.8cm}}
  \par\vspace{0.2em}
  Ahora verifica que la distribución es válida.}{Predice primero}
\renglones[4]
\respuesta{Es válida: las probabilidades suman 1. El promedio es 2,5 (se
calcula en el ejercicio 4).}
\solucion{%
  $0{,}10+0{,}15+0{,}20+0{,}25+0{,}30 = 1$, y ninguna es negativa. Es válida.
  \par
  Que sumen 1 no es un detalle: es el segundo axioma. Si una tabla de
  probabilidades no suma 1, está mal y no hay nada más que revisar.
  \par
  El promedio da $2{,}5$, y se calcula en el ejercicio 4. La estimación más
  frecuente es 3, porque uno mira dónde está la probabilidad más alta — y la
  moda no es la media.}

\ejercicio{1}{Con la misma variable, calcula la probabilidad de que un paciente
  asista a \textbf{al menos} 3 sesiones, y la de que abandone antes de la
  primera.}{}
\renglones[3]
\respuesta{$P(X\ge 3)=0{,}55$; $P(X=0)=0{,}10$.}
\solucion{%
  $P(X\ge 3)=P(3)+P(4)=0{,}25+0{,}30=0{,}55$.
  \par
  $P(X=0)=0{,}10$: uno de cada diez pacientes no llega a la primera sesión.
  \par
  Las dos son sumas de casillas de la tabla. Todavía no hace falta ninguna
  fórmula nueva — una distribución es sólo una lista de probabilidades bien
  organizada.}

\ejercicio{2}{Dos preguntas sobre la misma tabla que suenan parecidas y piden
  cosas distintas. Contesta las dos y explica qué cambia.
  \begin{enumerate}[label=(\alph*),topsep=2pt,itemsep=1pt]
    \item ¿Cuál es el número de sesiones \emph{más probable}?
    \item ¿Cuál es el número de sesiones \emph{esperado}?
  \end{enumerate}}{Dos casos casi iguales}
\renglones[5]
\respuesta{(a) 4 sesiones (la moda). (b) $E(X)=2{,}5$ (la media).}
\solucion{%
  (a) El más probable es \textbf{4}, con $P=0{,}30$. Es la \emph{moda}: se lee
  buscando la probabilidad más alta.
  \par
  (b) El esperado es:
  \[ E(X)=0(0{,}10)+1(0{,}15)+2(0{,}20)+3(0{,}25)+4(0{,}30)=2{,}5 \]
  \par
  \textbf{Son distintos y los dos son correctos}, porque responden preguntas
  diferentes. Si tienes que apostar a un solo paciente, apuesta a 4. Si tienes
  que estimar cuántas sesiones vas a dar en total a 100 pacientes, usa 2,5 y te
  da 250.
  \par
  Confundirlas hace planificar de más: con la moda dirías 400 sesiones.}

\ejercicio{2}{Una coordinadora calcula así la varianza de esa variable:
  \begin{quote}\itshape
  $\operatorname{Var}(X) = \dfrac{(0-2{,}5)^2+(1-2{,}5)^2+(2-2{,}5)^2+(3-2{,}5)^2+(4-2{,}5)^2}{5} = 2{,}5$
  \end{quote}
  \textbf{Señala el paso exacto} donde se equivoca y explica por qué.}{Encuentra el error}
\renglones[5]
\respuesta{Promedió los cuadrados por igual en vez de ponderarlos por su
probabilidad. Lo correcto es $1{,}75$.}
\solucion{%
  El error está en \textbf{dividir por 5}. Eso trata a los cinco valores como
  si fueran igualmente probables, y no lo son: el 4 pesa el triple que el 0.
  \par
  Cada cuadrado hay que multiplicarlo por \emph{su} probabilidad:
  \[ \operatorname{Var}(X)=\sum (x_i-\mu)^2 P(x_i) \]
  \[ =6{,}25(0{,}10)+2{,}25(0{,}15)+0{,}25(0{,}20)+0{,}25(0{,}25)+2{,}25(0{,}30) \]
  \[ = 0{,}625+0{,}3375+0{,}05+0{,}0625+0{,}675 = 1{,}75 \]
  \par
  Y $\sigma=\sqrt{1{,}75}\approx 1{,}32$ sesiones.
  \par
  \textbf{Cómo detectarlo sin la respuesta}: si el procedimiento divide por la
  cantidad de valores en vez de usar las probabilidades, la tabla sobraba — y
  si la tabla sobra, algo está mal.}
\enclase{Pregunta por qué se divide por 5 en la varianza de una muestra y no
aquí. La respuesta es que allá cada dato vale $1/n$; aquí cada valor ya trae su
peso.}

\ejercicio{3}{En las 200 fichas del servicio, el puntaje de depresión tiene
  $\mu = 6{,}32$ y $\sigma = 4{,}64$.
  \begin{enumerate}[label=(\alph*),topsep=2pt,itemsep=1pt]
    \item ¿Qué significa cada número, en palabras y en puntos del cuestionario?
    \item ¿Entre qué valores cae la mayoría de las fichas, aproximadamente?
    \item El corte está en 10. ¿Queda por encima o por debajo de
      $\mu+\sigma$?
  \end{enumerate}}{}
\renglones[6]
\respuesta{(a) Promedio 6,32 puntos, dispersión típica 4,64. (b) Entre 1,68 y
10,96. (c) $\mu+\sigma=10{,}96$: el corte queda por debajo.}
\solucion{%
  (a) El puntaje promedio del archivo es $6{,}32$ puntos. La desviación dice
  que lo típico es apartarse unos $4{,}64$ puntos de ese promedio.
  \par
  (b) Entre $\mu-\sigma = 1{,}68$ y $\mu+\sigma = 10{,}96$. Si la distribución
  tiene forma de campana, ahí cae alrededor del 68\,\% de las fichas — la regla
  que se usa en 2.9.
  \par
  (c) $\mu+\sigma = 10{,}96$, así que el corte de 10 queda \textbf{por debajo}:
  no llega a una desviación estándar por encima del promedio.
  \par
  Eso ya anticipa algo: el corte no es un valor extremo, así que va a marcar a
  bastante gente. Cuánta exactamente se calcula en 2.9.}

\ejercicio{3}{El servicio va a tamizar 20 estudiantes nuevos y quiere saber
  cuántos positivos esperar, para reservar horas de entrevista.
  \textbf{¿Alcanza con saber que $\mu=6{,}32$ y $\sigma=4{,}64$?} Si te parece
  que sí, calcula; si no, di qué hace falta.}{¿Alcanzan los datos?}
\renglones[5]
\respuesta{No alcanza con $\mu$ y $\sigma$ solos: hace falta la proporción de
positivos, $p=0{,}215$. Con ella, $E(X)=20\times 0{,}215=4{,}3$.}
\solucion{%
  \textbf{Con $\mu$ y $\sigma$ solos no alcanza.} Esos dos números describen el
  \emph{puntaje}, y la pregunta es sobre \emph{cuántas personas} superan un
  corte: son dos variables distintas.
  \par
  Lo que hace falta es la proporción de positivos, $p=\dfrac{43}{200}=0{,}215$.
  Con ella:
  \[ E(X)=n\cdot p = 20\times 0{,}215 = 4{,}3 \]
  Se esperan unos 4 o 5 positivos.
  \par
  \textit{Con un supuesto adicional sí se puede llegar desde $\mu$ y $\sigma$:
  si la distribución fuera normal, se obtiene $p\approx 0{,}214$ — y eso es
  exactamente lo que se hace en 2.9. Pero es un supuesto, no un dato, y
  conviene saber cuándo lo estás usando.}}

\ejercicio{4}{Un estudiante objeta:
  \begin{quote}\itshape
  «$E(X)=2{,}5$ sesiones no puede estar bien: nadie asiste a media sesión.»
  \end{quote}
  \begin{enumerate}[label=(\alph*),topsep=2pt,itemsep=1pt]
    \item ¿Tiene razón? Explica qué significa realmente el valor esperado.
    \item Da otro ejemplo de psicología donde $E(X)$ tampoco sea un valor
      posible.
    \item ¿Para qué sirve entonces, si nadie va a obtener ese valor?
  \end{enumerate}}{}
\renglones[7]
\respuesta{(a) No tiene razón: es el promedio a la larga, no un resultado
posible. (b) Ej.: 4,3 positivos entre 20 estudiantes. (c) Para planificar
totales.}
\solucion{%
  (a) \textbf{No tiene razón}, aunque la observación es buena. $E(X)$ no es «el
  resultado más probable» ni «un resultado posible»: es el \textbf{promedio al
  que tienden los resultados si el experimento se repite muchas veces}.
  \par
  (b) Los $4{,}3$ positivos esperados entre 20 estudiantes. Nunca vas a
  encontrar 4,3 personas — vas a encontrar 3, 4, 5 o 6. Pero si repites el
  tamizaje muchas semanas, el promedio se acerca a 4,3.
  \par
  (c) \textbf{Para planificar totales, no casos individuales.} Con 100
  pacientes y $E(X)=2{,}5$, vas a dar unas 250 sesiones. Ese número sí es real
  y sí se usa: define cuántas horas contratar.
  \par
  Es el mismo malentendido que «la familia promedio tiene 1,9 hijos». Nadie
  tiene 1,9 hijos, y sin embargo el número informa cuántas escuelas hacen
  falta.}

\trampa{usar la moda para planificar porque es «el valor más probable»}%
  {suena razonable: si 4 es lo más frecuente, parece el mejor pronóstico}%
  {para un caso individual puede servir, pero para totales hay que usar la
   media. Con la moda, el programa de 100 pacientes estimaría 400 sesiones en
   vez de 250 — un 60\,\% de más}

\cierre{%
  Una variable aleatoria es una lista de valores con sus probabilidades, y dos
  números la resumen: dónde se agrupa y cuánto se dispersa.
  \par
  Lo que viene es el paso siguiente: en vez de armar la tabla contando casos,
  hay procesos tan frecuentes que su tabla ya está calculada. Tres de ellos
  cubren casi todo lo que un servicio necesita planificar.}
